
Este capítulo apresenta o contexto em que o trabalho se insere, a motivação para a comparação entre PostgreSQL com JSONB e MongoDB, os objetivos da pesquisa e a forma como o texto está organizado.

\section{Contextualização e Motivação}
O crescimento do volume de dados semiestruturados trouxe novos desafios
aos Sistemas de Gerenciamento de Bancos de Dados (SGBDs) relacionais, o que leva a
adoção de SGBDs NoSQL \cite{Eldahshan2022}. Entre eles, os orientados a documentos,
como o MongoDB\footnote{\url{https://www.mongodb.com/docs/manual/}}, ganharam popularidade com o uso do formato JSON em aplicações web, 
pois armazenam dados semiestruturados de forma nativa, com esquema flexível e permitindo que documentos de uma mesma coleção tenham 
atributos distintos. 
No entanto, carecem de recursos dos sistemas relacionais \cite{Chasseur2013}, como junções eficientes, 
chaves estrangeiras para manter a consistência entre registros relacionados e transações ACID (atomicidade, consistência, isolamento e durabilidade) envolvendo múltiplas entidades.
Por sua vez, o PostgreSQL\footnote{\url{https://www.postgresql.org/docs/current/}} passou a incorporar suporte a dados
semiestruturados em formatos como JSON, com esforços para evoluir esse suporte
para representações binárias e indexáveis mais adequadas a esse modelo de dados
\cite{Bartunov2013}.

Essa discussão é relevante em catálogos de produtos, cujos dados costumam ser heterogêneos. 
Por exemplo, pode-se ter em um mesmo catálogo uma pizza que possui opções de borda e tamanhos, 
uma bebida que possui volume e informação de temperatura e uma sobremesa que possui informação nutricional específica, 
ou seja, cada tipo de produto tem suas próprias
particularidades. A representação horizontal
tradicional, com uma coluna por atributo, é inadequada \cite{Agrawal2001}, pois a maioria das
colunas ficaria vazia para a maioria dos itens, e cada nova categoria de produto exigiria
alterar o esquema da tabela. Criar uma tabela para cada tipo de produto também não é viável,
pois listar o cardápio completo passa a exigir a junção de várias tabelas, algo que seria muito custoso.
Uma alternativa é manter os atributos comuns em colunas relacionais e armazenar a parte variável em uma coluna JSONB\footnote{\url{https://www.postgresql.org/docs/current/datatype-json.html}}.
Outra alternativa é abandonar o modelo relacional e armazenar todo o catálogo em um SGBD orientado a documentos, como o MongoDB.


Esse é o cenário da equipe de Catálogo do iFood\footnote{\url{https://www.ifood.com.br/}}, empresa em que a autora realiza estágio.
O serviço responsável por esse catálogo armazena os itens em uma tabela do PostgreSQL
com 55 colunas, das quais 11 são do tipo JSONB, que armazenam objetos como preços,
grupos de opções, horários, classificações, entre outros. Os atributos comuns a todos os
itens permanecem em colunas relacionais, enquanto os que variam ficam em colunas JSONB, o que
permite lidar com a variabilidade sem abrir mão da integridade transacional do PostgreSQL.

O JSONB resolve o problema de modelagem, mas o principal gargalo passa a ser o desempenho.
Em sistemas relacionais, a ausência de um esquema fixo no JSONB resulta em consultas analíticas
lentas \cite{Durner2021}. No PostgreSQL, acessar um campo dentro de um JSONB exige operadores
e índices específicos, como o \textit{Generalized Inverted Index} (GIN)\footnote{\url{https://www.postgresql.org/docs/current/gin.html}}, enquanto no MongoDB o acesso a um campo do documento é
nativo. A pergunta que motiva este
trabalho é se o MongoDB seria uma alternativa mais adequada para o padrão de dados de
um catálogo como o do iFood, ou o PostgreSQL com JSONB e índices oferece desempenho
comparável.

Este trabalho é um recorte desse problema. Por se tratar de um sistema de produção que
atende milhões de transações diárias, o iFood não permite a replicação de sua infraestrutura
e de seus dados reais para fins de pesquisa acadêmica. Por isso, utiliza-se dados
gerados sinteticamente, em escala controlada de até 1 milhão de itens, e projeta o problema
em um cenário de dados estruturados e semiestruturados. Das 11 colunas
JSONB do sistema real, 3 foram selecionadas, cada uma representando um padrão estrutural observado nas 11 colunas: um conjunto de campos fixos (igual para todos os itens do catálogo),
um conjunto de campos que varia conforme a categoria do item, ou seja, o tipo de item vendido, como pizza, bebida ou sobremesa. Por fim,
um \textit{array} de estruturas aninhadas, em que cada elemento possui seus próprios campos.

\section{Objetivos}

O principal objetivo deste trabalho é comparar o desempenho do PostgreSQL, com a ferramenta
JSONB, e do MongoDB, identificando em quais cenários cada tecnologia apresenta melhores resultados.
Para isso, são avaliadas operações de leitura e de escrita típicas de um catálogo de
restaurantes. Nas leituras, consideram-se filtros de igualdade e exclusão em campos
escalares, a junção entre itens e restaurantes, a busca por intervalo em um campo numérico
do JSONB, a igualdade em um atributo aninhado e a contenção em um array de objetos aninhados.
Nas escritas, consideram-se a inserção condicional (\textit{upsert}), a exclusão lógica (\textit{soft delete}) e a atualização parcial de um campo JSONB.

Cada operação é avaliada com e sem índices, comparando estratégias de indexação genéricas e
específicas em cada sistema. Além disso, analisa-se a escalabilidade de cada tecnologia
conforme a seletividade das consultas, o volume de dados e o tamanho do lote de escrita.
Para garantir uma comparação justa, os experimentos são executados em um ambiente
controlado, com os mesmos dados sintéticos e limites de recursos.

\section{Organização do Trabalho}

Este trabalho de conclusão de curso estrutura-se em quatro capítulos, dos quais este é o Capítulo~\ref{chapter:introducao}. No Capítulo~\ref{chapter:bibliografia}, é apresentada a fundamentação teórica dos principais conceitos e tecnologias empregadas neste trabalho, como PostgreSQL, JSONB, MongoDB, Índices, entre outros, juntamente com os trabalhos relacionados a este estudo. Em seguida, no Capítulo~\ref{chapter:desenvolvimento}, descreve-se a metodologia adotada, detalhando-se a implementação das duas aplicações de catálogo comparadas e os resultados obtidos em todos os experimentos. Por fim, no Capítulo~\ref{chapter:conclusao}, são apresentadas as considerações finais sobre os resultados alcançados, as contribuições do trabalho e propostas para trabalhos futuros.